% Solutions/hints for ch15 exercises (assembled into Appendix C)

\begin{solution}{exr:ch15-derivation}
(a)~Write $g(\pos)=1/\rho(\pos)-1/\rho_0$, so $U_{\mathrm{rep}}=\tfrac12 k_{\mathrm{rep}}g^2$
inside the influence region. By the chain rule
$\nabla U_{\mathrm{rep}}=k_{\mathrm{rep}}\,g\,\nabla g$ and
$\nabla g=-\rho^{-2}\nabla\rho$, hence
$\vect{F}_{\mathrm{rep}}=-\nabla U_{\mathrm{rep}}=k_{\mathrm{rep}}\,(1/\rho-1/\rho_0)\,\rho^{-2}\,\nabla\rho$.
Outside the region both $U_{\mathrm{rep}}$ and its gradient are zero, and at
$\rho=\rho_0$ the inside expression is zero as well, so the force is continuous.
(b)~For a disc with centre $\vect{c}$ and radius $r$, $\rho=\norm{\pos-\vect{c}}-r$
and $\nabla\rho=(\pos-\vect{c})/\norm{\pos-\vect{c}}$. For a segment, $\rho$ is
the distance to the closest point $\vect{s}$ of the segment (\cref{ch:ch02}) and
$\nabla\rho=(\pos-\vect{s})/\norm{\pos-\vect{s}}$; the gradient is continuous
everywhere except on the set of points that have two closest points, which the
drone crosses in an instant.
(c)~With $d=\norm{\pos-\pos_{\mathrm{g}}}$ and $\nabla d=(\pos-\pos_{\mathrm{g}})/d$,
the product rule gives
$\nabla(U_{\mathrm{rep}}d^n)=d^n\nabla U_{\mathrm{rep}}+n\,d^{n-1}U_{\mathrm{rep}}\nabla d$,
so
$\vect{F}'_{\mathrm{rep}}=d^n\vect{F}_{\mathrm{rep}}-\tfrac{n}{2}k_{\mathrm{rep}}(1/\rho-1/\rho_0)^2d^{\,n-1}\,(\pos-\pos_{\mathrm{g}})/d$.
The second term points towards the goal. At the goal $d=0$ and, for $n\ge2$,
both terms vanish; for $n=1$ the second term has the magnitude
$\tfrac12 k_{\mathrm{rep}}(1/\rho-1/\rho_0)^2$ in every direction of approach, so
the force is discontinuous at the goal (a conic pull).
\end{solution}

\begin{solution}{exr:ch15-equilibrium}
On the axis $y=4$ with $1<x<3$ the force is horizontal by symmetry, with
$F_x(x)=(8-x)-\bigl(\tfrac{1}{3-x}-\tfrac12\bigr)\tfrac{1}{(3-x)^2}$. It is
$7$ at $x=1$ (the edge of the influence region) and tends to $-\infty$ as
$x\to3$, so it has a root; bisection gives $x^*=2.4872$, which is where the
simulation of \cref{sec:ch15-localmin} stops. The Hessian of $U$ there has the
eigenvalues $-2.65$ (in $y$) and $36.97$ (in $x$), computed with central
differences of the analytic force: the equilibrium is a saddle, because the
repulsion pushes a slightly off-axis drone further off the axis faster than the
attraction pulls it back. Only the measure-zero set of starts exactly on the axis
is trapped. For the peanut-shaped obstacle of the basin experiment the stopping
point $(2.676,4.1)$ has the eigenvalues $9.87$ and $24.89$: a genuine local
minimum, whose basin contains $179$ of the $441$ starts.
\end{solution}

\begin{solution}{exr:ch15-stability}
(a)~Without obstacles and without a speed limit the update reads
$\pos_{k+1}=\pos_k-\dt\,k_{\mathrm{att}}(\pos_k-\pos_{\mathrm{g}})$, so the error
$\vect{e}_k=\pos_k-\pos_{\mathrm{g}}$ obeys $\vect{e}_{k+1}=(1-k_{\mathrm{att}}\dt)\vect{e}_k$.
The factor has modulus below one, and the error shrinks, exactly when
$0<k_{\mathrm{att}}\dt<2$; for $k_{\mathrm{att}}\dt<1$ the approach is monotone,
for $1<k_{\mathrm{att}}\dt<2$ the sign alternates (overshoot), and for
$k_{\mathrm{att}}\dt>2$ the error grows. The self-test of
\code{ch15\_potential\_fields.py} checks the cases $0.5$ and $1.5$.
(b)~Near any minimum $\pos^*$ of a smooth $U$, $\vect{F}(\pos)\approx-\mat{H}(\pos-\pos^*)$
with the Hessian $\mat{H}$, so the error is multiplied by $\mat{I}-\dt\mat{H}$;
along an eigenvector with eigenvalue $\lambda$ the factor is $1-\dt\lambda$, and
the condition becomes $\dt\,\lambda_{\max}<2$.
(c)~In the passage of width $1.0$ the stiffness at the midpoint is $81.0$, so
$\dt=0.01$ gives $0.81$ (stable, monotone) and $\dt=0.05$ gives $4.05$
(unstable); the narrower passages have the stiffnesses $204$ (width $0.8$) and
$668$ (width $0.6$), which grow roughly like the inverse fourth power of the width.
\end{solution}

\begin{solution}{exr:ch15-gnron}
Hint: with the factor $d^n$ the total potential is
$U'=\tfrac12 k_{\mathrm{att}}d^2+U_{\mathrm{rep}}d^n\ge0$ and $U'(\pos_{\mathrm{g}})=0$,
so the goal is the global minimum; with $n\ge2$ the force there is zero by the
formula of \cref{exr:ch15-derivation}(c). The self-test checks both facts for the
goal $(5.4,4)$. For $n=1$ the goal is still the global minimum but the force does
not vanish at it; the drone reaches the goal from every side along a straight
final approach and would chatter around it without a goal tolerance. To see that
the correction does not remove local minima elsewhere, keep $n=2$ and rerun the
start $(0,4)$ of \cref{sec:ch15-localmin}: the drone still stops on the axis, at a
slightly different $x$.
\end{solution}
